% !TeX root = ../tg.tex

\chapter{绪论}
\section{课题的研究背景与意义}
随着人工智能与移动互联网技术的迅猛发展，分布式学习已演变为一种关键范式：它能在严格保障隐私安全的同时，充分释放海量私有数据的潜在价值。在这一领域中，联邦学习（Federated Learning, FL）采用一种协同训练范式：客户端仅利用本地数据更新模型，服务器则负责周期性聚合各客户端的模型参数以生成全局模型。这种‘共享模型参数而非原始数据’的机制，有效规避了数据隐私泄露风险。然而，传统联邦学习要求客户端在每轮训练中完整加载并更新上一轮的全局模型。面对日益庞大的深度学习模型，这一机制对客户端的计算算力与内存资源提出了严峻挑战，致使该方案在资源受限设备上的部署效率低下，甚至难以落地。

为突破这一瓶颈，分割联邦学习（Split Federated Learning, SFL）应运而生。SFL 借鉴了分割学习（Split Learning, SL）的模型切分思想，将神经网络模型划分为客户端模型和服务器端模型。在训练过程中，客户端仅负责前几层的计算，剩余层由服务器完成。这种协作机制不仅继承了联邦学习的隐私保护特性（客户端数据特征仅对本地可见），还显著降低了客户端的计算负载，使得在资源受限环境下部署大规模模型成为可能。尽管 SFL 在理论上兼具隐私与效率优势，但其在实际落地应用中，仍面临着数据异构性与系统资源受限的双重挑战，严重制约了其性能与普及。

首先，在数据层面，SFL 的性能高度依赖于客户端数据集的质量，而现实场景中的数据往往具有显著的异构性与噪声。 一方面，客户端数据通常呈现非独立同分布（Non-IID）特性，这种数据分布的差异性使得分布式训练过程更加脆弱。另一方面，数据收集或处理过程中难免产生标签噪声（Label Noise），即样本标签与真实值不一致。这种噪声可能源于人为标注错误、设备故障或传感器失效等多种原因。现有联邦学习研究表明，标签噪声会显著降低模型的鲁棒性、收敛速度及最终准确率。在 SFL 架构下，Non-IID 数据与标签噪声的耦合作用可能进一步放大训练的不稳定性，而目前针对 SFL 在复杂异构数据（特别是含噪 Non-IID 数据）下的收敛性理论分析与算法优化研究尚不充分。

其次，在系统资源层面，尽管 SFL 通过模型切分减轻了客户端负担，但其训练过程仍依赖标准的反向传播（Backpropagation, BP）算法，这带来了额外的资源开销。 标准 BP 方法要求客户端在计算过程中维护所有层的活动值（Activations）并计算精确的模型梯度，这导致了较高的内存消耗与计算成本。此外，SFL 中客户端与服务器之间需要在每个小批量（Mini-batch）数据上交换活动值与梯度，引入了显著的通信开销。对于电池容量有限、算力不足且网络带宽受限的边缘设备而言，现有的 SFL 训练机制仍显沉重，亟需设计更高效的算法以进一步节约通信、内存及计算资源。

综上所述，虽然分割联邦学习为隐私保护与资源受限场景下的分布式训练提供了新路径，但其在异构数据环境下的收敛稳定性以及在极端资源约束下的执行效率仍是亟待解决的关键科学问题。因此，本论文 titled“面向异构数据与资源受限场景的分割 - 联邦学习协同优化研究”，旨在从数据与系统两个维度出发，分别针对含噪异构数据的收敛性证明与算法设计，以及资源受限场景下的轻量化算法优化开展深入研究，以期实现分割联邦学习在复杂现实环境中的高效、稳健协同优化。
\section{国内外研究现状及分析}
\section{研究内容}
本研究围绕分割-联邦学习展开，聚焦于两个核心问题：异构数据对分割-联邦学习性能的影响及优化方案 ，以及轻量化的客户端训练机制设计与实现 。具体研究内容包括以下三个方面：
\subsection{收敛性分析与算法设计：异构数据对分割-联邦学习的影响以及优化方案}
\subsubsection{分割-联邦学习的收敛性分析}
在分割-联邦学习框架下，模型被纵向划分并分别部署于客户端与服务器端，客户端仅负责浅层网络的计算，而深层网络由服务器统一维护与更新。这种架构虽降低了客户端的计算负担，但也引入了新的挑战，尤其是在面对非独立同分布（Non-IID）数据和标签噪声时，其收敛行为尚不明确。

本研究系统地分析了分割-联邦学习在均方误差（MSE）和交叉熵损失函数下的收敛特性。通过建立理论框架，推导了在存在标签噪声的情况下，客户端本地随机梯度的方差上界及其与期望梯度之间的偏差。在此基础上，进一步探讨了不同数据分布、模型拆分点、本地训练轮次等因素对全局模型收敛速度和最终性能的影响。实验验证表明，在合理配置训练参数的前提下，分割-联邦学习可以在标签噪声和非独立同分布数据下稳定收敛，并保持较高的模型精度。

该部分的研究为后续鲁棒性算法的设计提供了理论支撑，也为分割-联邦学习在资源受限设备上的应用奠定了基础。
\subsubsection{对标签噪声和非独立同分布数据具有鲁棒性的分割-联邦学习算法}
由于实际应用场景中数据往往具有异构性和噪声特征，如何提升分割-联邦学习对此类数据的鲁棒性成为关键问题。为此，本研究提出了 SFL-DCT 算法，旨在有效识别并修正标签噪声，同时缓解非独立同分布数据带来的负面影响。

SFL-DCT 分为三个阶段：第一阶段通过局部迭代训练提取每个客户端的数据特征，并基于 LID（Local Intrinsic Dimensionality）或 LE（Label Entropy）指标检测可能含有噪声标签的客户端；第二阶段利用干净客户端集合训练一个鲁棒模型，并用于识别噪声样本并进行标签修正；第三阶段则在修正后的数据集上进行标准的分割-联邦学习训练。整个过程结合高斯混合模型（GMM）分类技术，实现了对噪声客户端和噪声样本的自动识别与处理。

实验结果表明，SFL-DCT 在 CIFAR-10、CIFAR-100 和 Fashion-MNIST 等多个数据集上均表现出良好的鲁棒性。即使在高达 60\% 的标签噪声率下，SFL-DCT 仍能维持较高的测试准确率，并优于传统联邦学习中的噪声鲁棒算法（如 FedCorr 和 FedNoRo）。此外，在 Non-IID 场景下，SFL-DCT 显著提升了模型的泛化能力，证明了其在异构环境中的有效性。
\subsection{轻量化客户端训练方案：无反向传播训练在分割-联邦学习中的应用}
为了进一步降低客户端的计算和内存开销，本研究首次将无反向传播训练机制引入分割-联邦学习框架中。传统的分割-联邦学习客户端需要执行前向传播和部分反向传播操作，这在资源受限设备上仍具有一定负担。为此，我们探索了多种无需依赖完整梯度反传的训练方法，如直接反馈对齐（Direct Feedback Alignment）、固定反馈权重等策略，并将其适配到分割-联邦学习的训练流程中。

具体而言，客户端侧采用无反向传播机制完成局部模型更新，仅需前向计算与少量反馈信号即可调整参数。这种设计显著减少了内存占用和计算复杂度，同时降低了通信成本。实验结果显示，相较于基于标准 SGD 的分割-联邦学习方法，引入无反向传播机制后，客户端的训练时间平均减少约 30\%，且在多个视觉任务上的模型精度下降幅度控制在 2\% 以内，展现了良好的性能与效率平衡。

该部分工作不仅拓展了分割-联邦学习的适用场景，还为边缘设备上的低功耗训练提供了新思路，具有较强的工程实用价值。